\documentclass[12pt]{article}

\usepackage{sbc-template}
\usepackage{graphicx,url}
\usepackage[utf8]{inputenc}
\usepackage{booktabs}
\usepackage{adjustbox}
\usepackage{amsmath}
\usepackage{float}

\sloppy

\title{Engenharia de Software 3.0: Uma Análise Empírica da\\
Manutenibilidade de Contribuições de Agentes Autônomos\\
de IA em Repositórios Open Source no GitHub}

\author{Arthur Curi, Helio Ernesto,\\
        João Campos, Lucas Azevedo,\\
        Luiz Saud, Nicolás Chabot,\\
        Raphael Brito, Yan Cota}

\address{Instituto de Ciências Exatas e Informática (ICEI)\\
  Pontifícia Universidade Católica de Minas Gerais (PUC Minas)\\
  Belo Horizonte --- MG --- Brasil
  \email{arthur.curi@sga.pucminas.br, helio.hernesto@sga.pucminas.br,
         joao.campos@sga.pucminas.br, lucas.azevedo@sga.pucminas.br,
         luiz.saud@sga.pucminas.br, nicolas.chabot@sga.pucminas.br,
         raphael.brito@sga.pucminas.br, yan.cota@sga.pucminas.br}
}

\begin{document}

\maketitle

\begin{abstract}
This paper investigates the maintainability of contributions generated by
autonomous Artificial Intelligence agents in open source GitHub repositories,
comparing them with human contributions. Using the AIDev-pop corpus and the
Goal-Question-Metric (GQM) approach, the study analyzes three dimensions:
structural code maintainability, post-merge indicators, and structural
technical debt. From 252{,}215 classified PRs, the analysis dataset comprises
202{,}014 PRs: only 9 in the AI Group (G1) and 202{,}005 in the Human Group
(G2); 50{,}201 ambiguous PRs (G3) were excluded. \textbf{None of the seven
research-question hypotheses were rejected after correction for multiple
tests} (Benjamini-Hochberg and Bonferroni). The only variable showing a
robust difference was the language distribution (Chi-square, adjusted
$p = 0{,}0236$), a descriptive control with negligible effect size. Raw
signals --- higher cyclomatic complexity, lower Maintainability Index, and
fewer test files in G1 --- did not survive correction and rest on a very small
G1 sample (n = 9; n = 3 for static code metrics), which severely limits
statistical power and external validity.
\end{abstract}

\begin{resumo}
Este artigo investiga a manutenibilidade de contribuições geradas por agentes
autônomos de Inteligência Artificial em repositórios open source no GitHub,
comparando-as com contribuições humanas. Utilizando o corpus AIDev-pop e a
abordagem Goal-Question-Metric (GQM), o estudo analisa três dimensões:
manutenibilidade estrutural do código, indicadores pós-merge e dívida técnica
estrutural. De 252.215 PRs classificados, o dataset de análise compreende
202.014 PRs: apenas 9 no Grupo IA (G1) e 202.005 no Grupo Humano (G2); 50.201
PRs ambíguos (G3) foram excluídos. \textbf{Nenhuma das sete hipóteses das
questões de pesquisa foi rejeitada após correção para múltiplos testes}
(Benjamini-Hochberg e Bonferroni). A única variável com diferença robusta foi
a distribuição de linguagens (Qui-quadrado, $p$ ajustado $= 0{,}0236$), um
controle descritivo com tamanho de efeito negligível. Sinais brutos ---
complexidade ciclomática mais alta, Maintainability Index mais baixo e menor
proporção de arquivos de teste em G1 --- não sobreviveram à correção e
baseiam-se em uma amostra G1 muito pequena (n = 9; n = 3 para métricas
estáticas de código), o que limita severamente o poder estatístico e a
validade externa.
\end{resumo}


% ─────────────────────────────────────────────────────────────────
\section{Introdução}

A adoção de agentes autônomos de Inteligência Artificial no desenvolvimento
de software tem crescido rapidamente nos últimos anos, configurando o que
vem sendo denominado de Engenharia de Software 3.0 \cite{li:25} ---
estágio em que sistemas baseados em modelos de linguagem passam a participar
ativamente de atividades como interpretar issues, escrever código, submeter
Pull Requests e interagir com revisores humanos \cite{vaithilingam:22,finnie:23}.
Esses agentes operam de forma cada vez mais autônoma, levantando questões
relevantes sobre o impacto que suas contribuições têm sobre a qualidade e a
sustentabilidade do código produzido.

Li, Zhang e Hassan \cite{li:25} mapearam essa transformação no corpus
AIDev, identificando 932.791 Pull Requests associados a agentes de IA
distribuídos em 116.211 repositórios públicos do GitHub. O subconjunto
AIDev-pop, restrito a repositórios com mais de 100 estrelas, compreende
33.596 PRs agentic de 2.807 repositórios populares, produzidos pelos
agentes Devin, OpenAI Codex e Claude Code.

Apesar do volume crescente de contribuições agentic, grande parte das
evidências disponíveis sobre a qualidade desse código foi obtida em tarefas
sintéticas, como o benchmark SWE-bench \cite{jimenez:24}, que medem
capacidade funcional em problemas isolados. Falta de estudos empíricos em
Pull Requests reais de projetos open source dificulta avaliar se contribuições
agentic são estruturalmente sustentáveis e qual o custo de manutenção que
podem introduzir para desenvolvedores humanos responsáveis pela evolução
do sistema \cite{dakhel:23}.

Este estudo investiga, com base em PRs reais, se as contribuições de agentes
autônomos de IA diferem das contribuições humanas em três dimensões de
manutenibilidade definidas pela ISO/IEC 25010: manutenibilidade estrutural
do código, indicadores pós-merge associados ao esforço de manutenção e
dívida técnica estrutural. A análise emprega métricas de análise estática
(complexidade ciclomática, code smells, Maintainability Index) e indicadores
de processo (taxa de reversão, número de revisões, presença de revisão),
estruturadas conforme a abordagem Goal-Question-Metric (GQM)
\cite{basili:94}. Os agentes considerados são os classificados como
autônomos pelo AIDev-pop --- Devin, OpenAI Codex e Claude Code --- em
contribuições escritas em JavaScript, TypeScript e Python.

O restante do artigo está organizado da seguinte forma. A Seção
\ref{sec:related} apresenta os trabalhos relacionados. A Seção
\ref{sec:method} descreve a metodologia, incluindo a base de dados,
o procedimento de identificação de PRs agentic, as questões de pesquisa
com hipóteses formais, a extração de métricas, o pipeline metodológico e
a análise estatística. A Seção \ref{sec:results} apresenta os resultados
por questão de pesquisa, encerrando com uma síntese dos testes e uma
discussão das limitações e ameaças à validade.


% ─────────────────────────────────────────────────────────────────
\section{Trabalhos Relacionados}
\label{sec:related}

\textbf{AIDev-pop e o paradigma de Engenharia de Software 3.0.} Li, Zhang
e Hassan \cite{li:25} apresentaram o corpus AIDev, que cataloga 932.791
PRs agentic em 116.211 repositórios do GitHub, identificando contribuições
de cinco agentes (Devin, OpenAI Codex, GitHub Copilot, Cursor e Claude
Code) por meio de marcadores explícitos em metadados de PR e padrões de
conta. O subconjunto AIDev-pop, restrito a repositórios com mais de 100
estrelas, abrange 33.596 PRs agentic de 2.807 repositórios populares. O
estudo original concentra-se em métricas de produtividade --- taxa de
merge, turnaround time e distribuição de linguagens --- sem investigar
atributos de manutenibilidade. Este trabalho é metodologicamente
complementar: utiliza o mesmo corpus para analisar qualidade estrutural,
indicadores pós-merge e dívida técnica.

\textbf{Análise de código gerado por Copilot.} Dakhel et al.
\cite{dakhel:23} avaliaram a qualidade do código produzido pelo GitHub
Copilot em tarefas experimentais, identificando problemas estruturais
recorrentes, incluindo redundâncias, padrões inadequados e dificuldades de
generalização em problemas mais complexos. O estudo demonstra que, apesar
de o assistente acelerar a produção, o código gerado frequentemente exige
revisões e ajustes adicionais. A principal diferença em relação ao presente
estudo está no contexto observacional: enquanto Dakhel et al.\ analisaram
tarefas sintéticas controladas, este trabalho examina PRs reais submetidos
em projetos open source de produção, aumentando a validade externa dos
resultados.

\textbf{Vieses metodológicos em Mining Software Repositories.}
Kalliamvakou et al.\ \cite{kalliamvakou:14} documentaram vieses recorrentes
em estudos baseados em GitHub, mostrando que repositórios pouco ativos,
forks e projetos pessoais podem distorcer análises empíricas se não forem
adequadamente filtrados. Os autores recomendam critérios mínimos de
qualidade de repositório, como atividade recente, número de estrelas e
exclusão de forks. Este estudo aplica diretamente essas recomendações: o
AIDev-pop seleciona repositórios não arquivados, não-forks e com mais de
100 estrelas, mitigando vieses associados a projetos de baixa relevância.

\textbf{Fingerprinting comportamental de agentes de IA.} Ghaleb
\cite{ghaleb:26} propôs uma técnica de identificação de PRs agentic
baseada em 41 atributos comportamentais --- padrões de mensagem de commit,
estrutura do PR e características do código --- atingindo F1-score de
97,2\% na classificação multi-classe de cinco agentes. O autor demonstra
que datasets que distinguem agentes apenas pela identidade do submissor
podem conter dados mal rotulados, especialmente em fluxos híbridos em que
desenvolvedores submetem código gerado por IA sob suas próprias contas.
Este estudo incorpora fingerprinting como segunda camada de validação,
aumentando a confiança nos rótulos do AIDev-pop e reduzindo o risco de
contaminação entre os grupos.

\textbf{Code smells como indicadores de manutenibilidade.} Palomba et al.\
\cite{palomba:18} conduziram estudo de larga escala sobre a difusão e o
impacto de code smells em centenas de projetos open source,
estabelecendo evidência empírica de que classes afetadas por smells
apresentam maior probabilidade de defeitos e maior esforço de manutenção.
O trabalho fundamenta o uso de densidade de code smells como métrica
operacional de manutenibilidade alinhada à ISO/IEC 25010. Este estudo
aplica esse arcabouço metodológico ao contexto específico de comparação
entre contribuições agentic e humanas.


% ─────────────────────────────────────────────────────────────────
\section{Metodologia}
\label{sec:method}

A Figura~\ref{fig:pipeline} sintetiza o pipeline metodológico do estudo.
A base é o AIDev-pop \cite{li:25}, subconjunto do AIDev com 33.596 PRs
agentic de 2.807 repositórios populares (>100 estrelas, não arquivados,
não-forks) em JavaScript, TypeScript e Python (2023--2024). Os agentes
analisados são Devin, OpenAI Codex e Claude Code. Critérios de inclusão:
diff entre 10 e 500 linhas, status \texttt{merged} e janela temporal
ancorada na data dos primeiros PRs agentic de cada repositório.

A identificação de PRs agentic adota três camadas: (1) metadados
estruturados do corpus (\texttt{devin-ai-integration[bot]},
\texttt{head:copilot/}, \texttt{Co-Authored-By: Claude}) \cite{li:25};
(2) fingerprinting comportamental \cite{ghaleb:26} com F1 de 97,2\%; e
(3) heurísticas de rastreabilidade \cite{mens:26} sobre artefatos do
repositório (\texttt{CLAUDE.md}, \texttt{.cursor/rules}, padrões de branch).
PRs cuja classificação não atingiu confiança suficiente foram alocados ao
Grupo 3 (ambíguos) e excluídos da análise principal.

Após a classificação, foram obtidos 252.215 PRs: 9 no Grupo IA (G1),
202.005 no Grupo Humano (G2) e 50.201 no Grupo Ambíguo (G3, excluído).
O dataset de análise comparativa reúne, portanto, 202.014 PRs (G1 + G2).
As métricas de processo (\texttt{review\_count}, \texttt{has\_review},
\texttt{is\_revert}, \texttt{days\_to\_merge}, \texttt{test\_file\_ratio})
foram computadas para todos os 202.014 PRs. As métricas estáticas que
exigem download e parsing do diff (CCM, Maintainability Index) foram
extraídas em uma amostra estratificada --- todos os PRs de G1 mais uma
amostra de G2 ---, resultando em n = 3 (G1) e n = 32 (G2) PRs com funções
analisáveis para CCM/MI, e n = 9 (G1) e n = 50 (G2) para densidade de smells
e tipos de smell.

O estudo investiga sete RQs em três dimensões, com hipóteses H0 (ausência
de diferença) e H1 (presença de diferença) e $\alpha = 0{,}05$.
\textbf{Q1 -- Estrutural}: RQ01 (CCM), RQ02 (smells/KLoC). \textbf{Q2 --
Pós-Merge}: RQ03 (reversão), RQ04 (rodadas de revisão). \textbf{Q3 --
Dívida técnica}: RQ05 (smells de design/KLoC), RQ06 (Maintainability
Index), RQ07 (tipos de smell).

CCM \cite{mccabe:76} e MI \cite{heitlager:07} são extraídos pelo Lizard
em modo incremental sobre o diff. Nesta execução, os detectores ESLint e
Pylint não foram acionados; a densidade de smells de design foi aproximada
pelo Lizard (funções longas, NLOC $>$ 50, normalizadas por KLoC), e o MI foi
calculado sem o volume de Halstead, o que tende a subestimar ligeiramente o
índice. Indicadores pós-merge vêm da GitHub API; LoC Delta é controle. Para
contínuas: Mann-Whitney U + Cliff's $\delta$ (Romano et al.: negligível
$|\delta|<0{,}147$; pequeno $<0{,}33$; médio $<0{,}474$; grande $\geq
0{,}474$). Binárias: Fisher; categóricas: Qui-quadrado. Todas ajustadas por
Benjamini-Hochberg e Bonferroni.

\begin{figure*}[!b]
\centering
\includegraphics[width=\textwidth]{metodologia.png}
\caption{Pipeline metodológico do estudo.}
\label{fig:pipeline}
\end{figure*}


% ─────────────────────────────────────────────────────────────────
\section{Resultados}
\label{sec:results}

A Tabela~\ref{tab:hasreview} apresenta a proporção de PRs com pelo menos
uma revisão humana formal.

\begin{table}[ht]
\centering
\caption{Proporção de PRs com Revisão Formal (\texttt{has\_review}).}
\label{tab:hasreview}
\begin{adjustbox}{max width=\textwidth}
\small
\begin{tabular}{lrrrl}
\toprule
\textbf{Grupo} & \textbf{N total} & \textbf{Com revisão (n)} & \textbf{Com revisão (\%)} & \textbf{$p$-value (Fisher)} \\
\midrule
IA (G1)        & 9        & 5        & 55,6\,\% & \multirow{2}{*}{$0{,}257$} \\
Humano (G2)    & 202.005  & 148.564  & 73,5\,\% & \\
\bottomrule
\end{tabular}
\end{adjustbox}
\end{table}

A diferença observada em \texttt{has\_review} (55,6\,\% em G1 contra
73,5\,\% em G2) \textbf{não é estatisticamente significativa}: o teste exato
de Fisher retornou $p = 0{,}257$. A seguir, apresentam-se os resultados por
questão de pesquisa.

\subsection{Q1: Manutenibilidade Estrutural}

\textbf{RQ01 -- CCM.} A Figura~\ref{fig:ccm} mostra a distribuição da
complexidade ciclomática média por grupo (n = 3 em G1; n = 32 em G2). Em
contraste com o esperado, G1 apresenta mediana mais alta ($\approx$ 5,4)
do que G2 ($\approx$ 1,25), com Cliff's $\delta = 0{,}71$ (efeito grande).
O teste Mann-Whitney U retornou $p$ bruto $= 0{,}043$ (limítrofe), mas
\textbf{a diferença não é estatisticamente significativa} após correção para
múltiplos testes ($p_{\text{BH}} = 0{,}142$; $p_{\text{Bonf}} = 0{,}681$).
Dado o tamanho amostral mínimo de G1 (n = 3), o resultado deve ser
interpretado como indício exploratório, não como evidência robusta.
\textbf{H0 não rejeitada} com $\alpha = 0{,}05$.

\begin{figure}[H]
\centering
\includegraphics[width=.55\textwidth]{violin_ccm_mean.png}
\caption{Distribuição de \texttt{ccm\_mean} por grupo (n = 3 G1; n = 32 G2).
G1 apresenta mediana mais alta ($\approx$ 5,4) que G2 ($\approx$ 1,25).
Diferença não significativa após correção; amostra G1 muito pequena.}
\label{fig:ccm}
\end{figure}

\textbf{RQ02 -- Densidade de Code Smells / KLoC.} Os detectores ESLint e
Pylint não foram executados nesta análise; o indicador disponível é a
densidade de smells de design aproximada pelo Lizard (ver RQ05). Sobre essa
proxy, o teste Mann-Whitney U não indicou diferença estatisticamente
significativa entre G1 e G2 ($p = 0{,}612$; $\delta = 0{,}049$, negligível).
\textbf{H0 não rejeitada} com $\alpha = 0{,}05$.

\subsection{Q2: Indicadores Pós-Merge}

\textbf{RQ03 -- Taxa de Reversão.} Nenhum dos 9 PRs de G1 foi revertido
(0\,\%), contra 6.813 de 202.005 em G2 ($\approx$ 3,4\,\%). O teste exato
de Fisher retornou $p = 1{,}0$, sem diferença estatisticamente significativa.
\textbf{H0 não rejeitada} com $\alpha = 0{,}05$.

\textbf{RQ04 -- Número de Rodadas de Revisão.} A Figura~\ref{fig:review}
apresenta a distribuição de \texttt{review\_count} por grupo. Ambos os grupos
têm mediana igual a 1; G1 concentra-se em 0--4 revisões, enquanto G2
apresenta cauda superior estendida até 30. O teste Mann-Whitney U retornou
$p = 0{,}411$ com Cliff's $\delta = -0{,}153$ (efeito pequeno), \textbf{sem
diferença estatisticamente significativa}. \textbf{H0 não rejeitada} com
$\alpha = 0{,}05$.

\begin{figure}[H]
\centering
\includegraphics[width=.55\textwidth]{violin_review_count.png}
\caption{Distribuição de \texttt{review\_count} por grupo
($p = 0{,}411$; Cliff's $\delta = -0{,}153$). Mediana igual a 1 em ambos;
diferença não significativa.}
\label{fig:review}
\end{figure}

A Figura~\ref{fig:days} mostra a distribuição de \texttt{days\_to\_merge}
em escala logarítmica. G1 concentra-se próximo a zero (mediana $\approx$
0,04 dia), assim como G2 (mediana $\approx$ 0,62 dia), que exibe maior
variância e outliers até $\sim$7,7 (em log). O teste Mann-Whitney U
\textbf{não indicou diferença significativa} ($p = 0{,}305$; $\delta =
-0{,}197$, pequeno).

\begin{figure}[H]
\centering
\includegraphics[width=.55\textwidth]{violin_days_to_merge.png}
\caption{Distribuição de \texttt{days\_to\_merge} em escala log
($p = 0{,}305$; $\delta = -0{,}197$). Diferença não significativa.}
\label{fig:days}
\end{figure}

\subsection{Q3: Dívida Técnica Estrutural}

\textbf{RQ05 -- Densidade de Code Smells de Design / KLoC.} O teste
Mann-Whitney U sobre \texttt{design\_smell\_density} (n = 9 G1; n = 50 G2;
proxy via Lizard) não indicou diferença estatisticamente significativa entre
os grupos ($p = 0{,}612$; $\delta = 0{,}049$, negligível). Ambas as medianas
são iguais a zero. \textbf{H0 não rejeitada} com $\alpha = 0{,}05$.

\textbf{RQ06 -- Maintainability Index.} A Figura~\ref{fig:mi} apresenta
a distribuição de \texttt{maintainability\_index} por grupo (n = 3 G1; n = 32
G2). G1 concentra-se em torno de 49--50 (mediana $\approx$ 49,6); G2
apresenta mediana mais alta ($\approx$ 61,4), com Cliff's $\delta = -0{,}625$
(efeito grande). O teste Mann-Whitney U retornou $p = 0{,}082$, \textbf{sem
diferença significativa}. Apesar do tamanho de efeito grande, o resultado não
é robusto dado o n mínimo de G1 (n = 3). \textbf{H0 não rejeitada} com
$\alpha = 0{,}05$.

\begin{figure}[H]
\centering
\includegraphics[width=.55\textwidth]{violin_maintainability_index.png}
\caption{Distribuição de \texttt{maintainability\_index} por grupo (n = 3 G1;
n = 32 G2). G1 mediana $\approx$ 49,6; G2 mediana $\approx$ 61,4. Diferença
não significativa após correção; amostra G1 muito pequena.}
\label{fig:mi}
\end{figure}

\textbf{RQ07 -- Distribuição dos Tipos de Code Smell.} A
Figura~\ref{fig:smell} mostra a distribuição de \texttt{smell\_type} (n = 9
G1; n = 50 G2). Ambos os grupos são dominados por PRs sem smell
(\textit{none}: 78\,\% em G1, 76\,\% em G2). G1 apresenta maior proporção de
smells de \textit{complexity} (22\,\% vs.\ 6\,\%), ao passo que G2 concentra
os smells de \textit{style} (12\,\% vs.\ 0\,\%) e \textit{size} (6\,\% vs.\
0\,\%). O teste Qui-quadrado retornou $p = 0{,}262$, \textbf{sem diferença
estatisticamente significativa}. \textbf{H0 não rejeitada} com $\alpha =
0{,}05$.

\begin{figure}[H]
\centering
\includegraphics[width=.55\textwidth]{bar_smell_type.png}
\caption{Distribuição dos tipos de \texttt{smell\_type} por grupo (n = 9 G1;
n = 50 G2). Predominância de \textit{none} (78\,\% G1; 76\,\% G2). Diferença
não significativa.}
\label{fig:smell}
\end{figure}

\subsection{Síntese dos Testes}

A Tabela~\ref{tab:resumo} consolida os resultados estatísticos por métrica.
\textbf{Nenhuma das sete hipóteses das questões de pesquisa foi rejeitada
após correção para múltiplos testes.} A única variável com diferença robusta
foi a distribuição de linguagens (\texttt{language}), uma variável descritiva
de controle: o Qui-quadrado retornou $p = 0{,}0015$, mantendo significância
após Benjamini-Hochberg e Bonferroni ($p_{\text{adj}} = 0{,}0236$), embora com
tamanho de efeito negligível ($\delta = 0{,}008$). Esse resultado apenas
reflete a composição amostral --- G1 é majoritariamente Python (8 de 9 PRs),
enquanto G2 é dominado por TypeScript --- e não diz respeito à
manutenibilidade. Alguns sinais brutos (CCM mais alta e MI mais baixo em G1;
menor proporção de arquivos de teste e de arquivos alterados) surgiram com
$p < 0{,}05$ antes da correção, mas todos se dissolveram após o ajuste para
múltiplos testes e repousam sobre uma amostra G1 mínima.

\begin{table}[ht]
\centering
\caption{Resumo dos resultados estatísticos por métrica (estrato \texttt{all}).}
\label{tab:resumo}
\begin{adjustbox}{max width=\textwidth}
\small
\begin{tabular}{llllll}
\toprule
\textbf{Métrica} & \textbf{Teste} & \textbf{$p$ bruto} & \textbf{B-H} & \textbf{Bonf.} & \textbf{Cliff's $\delta$ / Decisão} \\
\midrule
\texttt{language}               & Qui-quadrado & $0{,}0015$ & Sig.     & Sig.     & $0{,}008$ (neglig.) -- controle, robusto \\
\texttt{ccm\_mean}              & Mann-Whitney & $0{,}043$  & Não sig. & Não sig. & $0{,}71$ (grande) -- H0 não rejeitada \\
\texttt{maintainability\_index} & Mann-Whitney & $0{,}082$  & Não sig. & Não sig. & $-0{,}625$ (grande) -- H0 não rejeitada \\
\texttt{review\_count}          & Mann-Whitney & $0{,}411$  & Não sig. & Não sig. & $-0{,}153$ (pequeno) -- H0 não rejeitada \\
\texttt{has\_review}            & Fisher       & $0{,}257$  & Não sig. & Não sig. & -- -- H0 não rejeitada \\
\texttt{design\_smell\_density} & Mann-Whitney & $0{,}612$  & Não sig. & Não sig. & $0{,}049$ (neglig.) -- H0 não rejeitada \\
\texttt{is\_revert}             & Fisher       & $1{,}000$  & Não sig. & Não sig. & -- -- H0 não rejeitada \\
\texttt{days\_to\_merge}        & Mann-Whitney & $0{,}305$  & Não sig. & Não sig. & $-0{,}197$ (pequeno) -- H0 não rejeitada \\
\texttt{smell\_type}            & Qui-quadrado & $0{,}262$  & Não sig. & Não sig. & H0 não rejeitada \\
\bottomrule
\end{tabular}
\end{adjustbox}
\end{table}

\subsection{Limitações e Ameaças à Validade}

O principal fator limitante deste estudo é o tamanho do Grupo IA. Dos 295 PRs
agentic identificados pelos rótulos do AIDev-pop nos repositórios coletados,
apenas 9 sobreviveram aos critérios de inclusão (diff de 10--500 linhas,
linguagens JS/TS/Python e janela temporal), e desses, somente 3 possuíam
funções analisáveis para as métricas estáticas de código (CCM e MI). Além
disso, os 9 PRs de G1 concentram-se em um único repositório
(\texttt{Significant-Gravitas/AutoGPT}), o que compromete a representatividade
e a validade externa. Com n = 9 (e n = 3 para CCM/MI), o poder estatístico é
muito baixo, e os efeitos de tamanho grande observados em \texttt{ccm\_mean}
e \texttt{maintainability\_index} não puderam ser confirmados após correção.
Os resultados devem, portanto, ser lidos como exploratórios e preliminares.

Há ainda limitações de instrumentação: os detectores ESLint e Pylint não
foram executados nesta rodada, de modo que a densidade de smells (RQ02 e RQ05)
foi aproximada por uma heurística do Lizard (funções longas por KLoC), não
capturando violações de convenção e de acoplamento; e o Maintainability Index
foi calculado sem o volume de Halstead, subestimando ligeiramente o índice.
A forte assimetria amostral (202.005 PRs humanos contra 9 de IA) também torna
os testes não paramétricos sensíveis a desbalanceamento. Trabalhos futuros
devem ampliar a amostra agentic --- relaxando os filtros de coleta ou focando
repositórios com maior concentração de PRs de agentes --- e habilitar os
detectores de smells completos antes de extrair conclusões sobre
manutenibilidade.


% ─────────────────────────────────────────────────────────────────
\bibliographystyle{sbc}
\bibliography{sbc-aiddev}

\end{document}
